\subsection{Principios de integración}
\label{subsec:principios-integracion}

La integración no es un apéndice: es el tejido que sostiene la operación distribuida de Puelche. El riesgo principal del caso está en las costuras entre sistemas legados sin documentación de interfaces, no en los módulos nuevos. De las Bases Técnicas Transversales (RT-02.06 a RT-02.08 y RT-05.16; Bases Técnicas Transversales, cap.~2, p.~7, y cap.~5, p.~12) y del caso 02 se desprenden ocho reglas de diseño:

\begin{enumerate}
  \item \textbf{Contrato antes que conexión.} Cada interfaz tiene dueño, contrato OpenAPI o AsyncAPI, versión y comportamiento ante falla antes de entrar en operación.
  \item \textbf{Asincronía con excepciones justificadas.} Los eventos desacoplan procesos; disponibilidad, crédito, autorización de pago y actos tributarios que exigen respuesta se consultan con tiempo máximo de espera y estado explícito.
  \item \textbf{Idempotencia.} Una escritura conserva su UUID en todos los reintentos; el consumidor deduplica y resuelve conflictos por regla de negocio, no por la última marca de tiempo.
  \item \textbf{Reconciliación auditable.} Cada conflicto conserva operación, regla, resultado, autor y momento, conforme a la bitácora del Artículo 16.4 (Bases Administrativas, art. 16.4, p. 12).
  \item \textbf{Frontera única del legado.} La capa anticorrupción traduce hacia el ERP; ningún módulo o portal escribe directamente en él.
  \item \textbf{Confianza explícita.} Cada mensaje y llamada se autentica, autoriza y correlaciona; las excepciones de conectividad hacia el sitio son privadas, acotadas y auditadas.
  \item \textbf{Falla declarada.} Las quince entradas del catálogo indican si reintentan, se difieren, degradan o requieren procedimiento manual.
  \item \textbf{Ventana de despacho protegida.} Las cargas masivas, reprocesos y despliegues de conectores no interrumpen la operación de 05:30 a 07:00.
\end{enumerate}

Estas reglas se verifican en los contratos, el catálogo de interfaces y las pruebas de falla de cada consumidor.

Estos principios se materializan en la capa de integración, los eventos canónicos, la capa anticorrupción, el catálogo de interfaces y las reglas de reconciliación descritas en este apartado.

El núcleo reúne M1--M12 en Laravel 13 sobre PHP 8.5, con dependencias fijadas por \texttt{composer.lock}. Cada módulo tiene espacio de nombres, servicios de aplicación, reglas de dominio y adaptadores propios. Esta organización evita que una pantalla o un cambio de proveedor invada otra capacidad. Para 14.200 clientes y unos 31.000 pedidos mensuales, se conserva la sencillez operativa del monolito modular, en línea con el numeral 2.3 de las Bases Técnicas Transversales. Sincronización, EDI y telemetría tienen trabajadores y capacidad de escalado independientes desde el inicio (RT-02.02; Bases Técnicas Transversales, cap. 2, p. 6); su emplazamiento se detalla en 4.2.

La volumetría operativa que condiciona el dimensionamiento es la siguiente: 14.200 clientes activos, 8.400 SKU (que pasan a aproximadamente 9.500), 31.000 pedidos mensuales con 260.000 líneas y 2,4 millones de unidades, aproximadamente 1.400 entregas diarias habituales con un pico de septiembre de 2.600 entregas (volumen casi duplicado durante tres semanas), 96 camiones en la ventana crítica de despacho (42 propios y 54 de transportistas externos), 62 preventistas, unos 160 conductores de transportistas y 310 personas de centro de distribución (Bases Técnicas del caso, cap.~14, pp.~24--25), además de 84 conductores propios y peonetas (Bases Técnicas del caso, cap.~2, p.~5). De esas 310 personas, 120 preparan pedidos en el turno nocturno de Talca (Bases Técnicas del caso, cap.~8, p.~14). El máximo horario global de septiembre alcanza 14,66 TPS a las 12:00, incluida la preventa y el portal; dentro de la ventana de despacho de 05:30 a 07:00, el máximo es 6,94 TPS. El diseño considera ambos picos según el proceso que dimensiona, nunca el promedio, y declara explícitamente los puntos únicos de falla (SPOF) con su mitigación conforme a RT-02.11.

El modelo distingue seis instalaciones: los CD de Talca y Concepción, las plataformas de cross-docking de Curicó, Chillán y Los Ángeles, y la casa matriz de Talca, contigua al CD principal (Bases Técnicas del caso, cap.~2, p.~5). Las primeras cinco ejecutan el mismo perfil \texttt{wms\_only}, conservan autoridad sobre los movimientos de su bodega y publican eventos idempotentes para consolidar stock y reservas en N-04/N-05; la casa matriz no aloja cómputo propio y utiliza los portales y las consolas en nube. Los identificadores de sitio serán parametrizables para incorporar una séptima instalación proyectada a tres años y evaluar la futura operación en Los Lagos hacia 2030.

\clearpage
